% Part: lambda-calculus
% Chapter: introduction
% Section: overview

\documentclass[../../../include/open-logic-section]{subfiles}

\begin{document}

\olfileid[fr]{lam}{int}{ovr}
\olsection{Vue d'ensemble}

Alonzo Church a conçu le lambda-calcul au début
des années 1930 comme fondement d'une logique
constructive, et \emph{non} comme modèle des
fonctions calculables. On a cependant vite montré
qu'il équivalait à d'autres définitions de la
calculabilité, notamment celles qui reposent sur
les fonctions calculables par machine de Turing
et sur les fonctions récursives partielles. Que
cette équivalence ait d'abord suscité une certaine
surprise rend la caractérisation d'autant plus
intéressante.

La notation lambda offre un moyen commode de
désigner une fonction par une expression
symbolique qui la définit, sans lui donner
d'abord un nom. Au lieu
de dire « Soit $f$ la fonction définie par
$f(x)=x+3$ », on peut dire « Soit $f$ la fonction
$\lambd[x][(x+3)]$ ». Autrement dit,
$\lambd[x][(x+3)]$ est simplement un \emph{nom}
de la fonction qui ajoute trois à son argument.
Dans cette expression, $x$ est une variable liée par l'abstraction.
On peut la renommer en~$y$ et désigner la même fonction par
$\lambd[y][(y+3)]$. La notation admet aussi des
paramètres supplémentaires. Supposons, par
exemple, que $g(x,y)$ soit une fonction de deux
variables et que $k$ soit un entier naturel.
Alors $\lambd[x][g(x,k)]$ est la fonction qui
envoie chaque~$x$ sur~$g(x,k)$.

Cette manière de définir une fonction à partir
d'une expression symbolique s'appelle
l'\emph{abstraction lambda}. Son opération
complémentaire est l'\emph{application} : si
l'on dispose d'une fonction~$f$ (définie, par
exemple, sur les entiers naturels), on peut
l'appliquer à une valeur telle que~$2$. Dans
la notation usuelle, le résultat s'écrit~$f(2)$.

L'application d'une abstraction se simplifie en substituant
l'argument aux occurrences de la variable liées par cette
abstraction. On renomme au besoin des variables liées dans
le corps pour éviter de capturer une variable libre de l'argument. Par exemple,
\[
(\lambd[x][(x+3)])(2)
\]
se simplifie en~$2+3$.

Jusqu'ici, nous avons seulement introduit de
nouvelles notations pour des notions usuelles.
Le lambda-calcul s'écarte cependant plus
radicalement du point de vue de la théorie des
ensembles. Dans ce cadre :
\begin{enumerate}
\item Tout désigne une fonction.
\item On peut définir des fonctions par
      abstraction lambda.
\item Tout terme peut être appliqué à tout autre terme.
\end{enumerate}
En particulier, pour tout terme~$F$ du lambda-calcul,
l'auto-application $F(F)$ est un terme bien formé. On
parle de \emph{lambda-calcul non typé}, puisque la
formation d'une application ne dépend d'aucun type.

\begin{digress}
Il existe aussi un \emph{lambda-calcul typé}, variante
importante du système non typé. Il lui ressemble à
bien des égards, mais s'accorde beaucoup plus
facilement avec le cadre classique de la théorie
des ensembles et possède certaines propriétés
très différentes.

Les recherches sur le lambda-calcul ont joué un
rôle central en informatique théorique et dans
la conception des langages de programmation.
LISP, conçu par John McCarthy dans les années
1950, est un exemple ancien de langage influencé
par ces idées.
\end{digress}

\end{document}
